% Generated draft: requires mathematical and source review.
\documentclass[11pt]{article}
\usepackage[margin=25mm]{geometry}
\usepackage{fontspec}
\setmainfont{Latin Modern Roman}
\usepackage{amsmath,amssymb,mathtools,unicode-math}
\setmathfont{Latin Modern Math}
\usepackage{enumitem,needspace,xurl,hyperref}
\setlength{\parindent}{0pt}
\setlength{\parskip}{5pt}
\setlength{\emergencystretch}{3em}
\newcommand{\N}{\mathbb{N}}
\newcommand{\Z}{\mathbb{Z}}
\newcommand{\Q}{\mathbb{Q}}
\newcommand{\R}{\mathbb{R}}
\newcommand{\C}{\mathbb{C}}
\newcommand{\F}{\mathbb{F}}
\newcommand{\A}{\mathbb{A}}
\newcommand{\PP}{\mathbb{P}}
\newcommand{\E}{\mathbb{E}}
\newcommand{\eps}{\varepsilon}
\newcommand{\dd}{\,\mathrm{d}}
\begin{document}
\section*{UnsolvedMath ID: 800894}
\textbf{Source problem code:} catalogue entry
\subsection{Definitions and mathematical statement}
The notation \(u\in\mathcal{C}^{2}(\Omega)\) expresses interior twice continuous differentiability: all partial derivatives of \(u\) of order at most two exist and are continuous on \(\Omega\). The supplied record does not specify the ambient space, its dimension, the domain \(\Omega\), the function \(u\), or the equation and boundary conditions that determine \(u\).

The source question is:
\begin{quote}
Under what minimal assumptions does \(u\) belong to \(\mathcal{C}^{2}(\Omega)\)?
\end{quote}
The intended hypotheses cannot be recovered from this excerpt: its phrase ``under the assumptions above'' refers to material that is not supplied. Thus this is an incompletely specified regularity question, rather than a fully quantified mathematical assertion. No particular equation, domain class, boundary data, or regularity hypotheses are imposed in this draft.

As background, the supplied excerpt attributes to Bedford and Fornaess, under the reference label [BeFo], a construction showing that \(u\) need not belong to \(\mathcal{C}^{2}(\bar\Omega)\) under those omitted assumptions. It says that their construction relies on the boundary behavior of the \(\mathcal{C}^{2}\) norm. This attribution is retained as source-reported background; its hypotheses and the precise convention for regularity on \(\bar\Omega\) are unavailable. The requested conclusion concerns the interior \(\Omega\), whereas the reported obstruction concerns regularity up to the boundary.
\subsection{Short English statement}
The question asks for minimal hypotheses ensuring that the intended function \(u\) is twice continuously differentiable in the interior of \(\Omega\). The defining problem and preceding assumptions are missing, so a precise quantified formulation requires additional source context.
\subsection{Sources}
\begin{itemize}
\item Local UnsolvedMath registry, supplied source record: permanent ID \(800894\), exact source code \texttt{catalogue entry}, title ``Open Problems in Pluripotential Theory --- Question 25.'' The registry attributes the question to Dinew, Guedj, and Zeriahi, with proposed year 2015. Its supplied status is ``open''; this label is not an independently verified assessment.
\item Dinew, Guedj, and Zeriahi, \emph{Open Problems in Pluripotential Theory} (2015), Question 25. Supplied source URL: \url{https://arxiv.org/abs/1511.00705}. The record reports access on 2026-10-16 and extraction by \texttt{arxiv\_latex}; these are supplied provenance details, not a claim of independent source access or verification.
\item Bedford and Fornaess, reference label [BeFo] in the supplied excerpt, credited there with the boundary regularity obstruction described above. No full bibliographic entry for [BeFo] is supplied.
\end{itemize}
\end{document}
